\section{Discussion and Research Outlook}\label{sec:9}

The SAGA comparison theorem is the first theorem about the
local-to-global capability of AAT to reach a full proof, and it is
also the first traversal of a larger research program. This section
discusses what the theorem has settled (\S\ref{sec:91}) and the
research directions it opens (\S\ref{sec:92}--\S\ref{sec:96}).
Everything from \S\ref{sec:92} onward is research outlook, not to be
confused with proved results; the development of theorems,
implementations, and tooling in each direction is outside the scope of
this paper.

\subsection{What SAGA has settled}\label{sec:91}

The first consequence of the comparison theorem is
\textbf{translatability}. A diagnosis posed in the language of
semantic repair and a computation posed in the language of equation
geometry agree at the level of $H^1$ and residual classes. Subsequent
theory development can proceed on either side and be translated to the
other: structure discovered on the semantic side can be computed with
geometric tools, and theorems proved on the geometric side can be read
in the language of repair.

The second consequence is the \textbf{grounding of counterfactuals}.
The sentence ``without this fix, the whole would not have glued'' used
to be a counterfactual with no place to stand within the vocabulary of
testing and observation, because a failure that did not occur cannot
be observed. The three-way equivalence of Theorem~\ref{thm:gluing}
gives this sentence a mathematical status. Since the zero class is
equivalent to the existence of a global repair, ``it glued because the
correction killed the class'' is a consequence, not a heuristic. The
before/after-repair comparison of Section~\ref{sec:7} is the first
working diagnostic example of how this consequence operates in an
engineering process (the division of labor between the scope of the
measurement checks and the Lean witnesses of Section~\ref{sec:6} is in
\S\ref{sec:77}).

\subsection{The Rising Sea as method}\label{sec:92}

The method of this research program follows what Grothendieck called
\emph{la mer qui monte} --- the rising sea \citep{grothendieck2022recoltes}.
Instead of prying open a hard problem with ad hoc tools, one raises
the water level of abstraction of the theory until the problem
dissolves naturally at the higher level.

The engineering problem ``locally correct but globally broken'' has,
at the level where one works directly with code and modules, remained
a target attacked with the local tools of testing, contracts, and CI.
Because the problem itself lives in the overlaps of parts, local
tools cannot reach it in principle. Rising to the level of Atoms,
sites, sheaves, covers, and cohomology, the same problem becomes an
ordinary object with coordinates: an end-to-end failure is a class in
$H^1$; classes can be computed; and what can be computed becomes the
work of tools.

This raising of the water level has two proper features. First,
\textbf{abstraction is portability}. Because Atoms are placed at a
level independent of language and framework, neither the theory nor
the measurement depends on a particular technology stack, and a
multi-language system can be treated as a single architecture
geometry. Second, \textbf{the water level is maintained by Lean}.
Grothendieck's sea settled in thousands of pages of prose; the sea of
this program is maintained as a tower of machine-verified theorems,
and as long as the theorems are formally verified, the water does not
recede. That the name SAGA is an homage to Serre's GAGA
\citep{serre1956gaga} --- the comparison theorem that established the
correspondence of two geometries --- reflects this configuration.

\subsection{The next mathematical peaks}\label{sec:93}

From the SAGA mathematics proved in this paper, the following theory
developments can be formulated naturally.
\begin{itemize}
\item \textbf{Characterization of the descent condition}: the
three-way equivalence of Theorem~\ref{thm:gluing} holds under the
selected true sheaf condition. The next peak is to characterize the
validity of this condition itself from architecture data and lower it
to checkable assumptions. Once realized, ``this architecture satisfies
the descent assumptions, so passing local verification implies global
consistency'' becomes a checkable measurement item, and the theory can
point to the places where the assumptions break as exactly the places
where integration verification should concentrate.
\item \textbf{Higher coherence}: developing as independent statements
what this paper deliberately separated (\S\ref{sec:57}): the
pointed-set-valued $H^1$ of nonabelian torsors, 2-cocycles and gerbes,
obstructions in $H^2$ and beyond, stack descent.
\item \textbf{Architecture schemes and morphism theory}: treating the
zero locus of the required equations as a scheme, and proving theorems
about morphisms, base change, and fibers between architectures.
\item \textbf{Moduli of repairs and derived deformation}: treating the
space of repairs itself as a geometric object, speaking of
deformations, degenerations, and singularities of repairs.
\item \textbf{First-order conormal specialization}: the specialization
of the conormal geometry given by the first-order approximation of the
obstruction ideal.
\item \textbf{Deepening the Atom foundations}: refinement of the Atom
axiom system and extension of the correspondence theory between
semantic atoms and equation Atoms.
\end{itemize}

\subsection{From statics to dynamics: Software Field Theory}\label{sec:94}

AAT is a statics: it treats the consistency of an architecture at a
point in time. Software Field Theory (SFT) is the research program
that constructs the dynamics of software evolution on top of that
geometry: organizing development traces as a site, treating changes,
branches, and merges in geometric language, and formulating review,
CI, and operational feedback as forces acting on the development
system.

The central viewpoint of SFT is a change in the role of architecture:
an architecture is not only the present shape of the code but the
\textbf{shape of the reachable futures}. A good architecture is a
field configuration in which desirable futures are reachable and
undesirable futures are hard to reach. Under this view, SFT aims at
central propositions of the following form:

\begin{quote}
A software architecture is modular exactly when its future evolution
satisfies descent.
\end{quote}

A merge in version control is a gluing of local changes, and SFT reads
it as descent theory. The static comparison proved by SAGA is the
ground this dynamics stands on at each instant. Where the static
$H^1$ measures ``does not glue now'', SFT aims to measure ``can this
sequence of changes reach a future that glues''.

\subsection{The future of measurement: a fixed point in the age of generation}\label{sec:95}

As research outlook, we also record one configuration on the
measurement side.

In development where code is generated faster than humans can
understand it, the code, the tests, the reviews, and the fix proposals
all sway on the same probabilistic generative ground. Within that, a
deterministic measurement grounded in theorems becomes a point of a
different nature: its verdict is not someone's opinion but the
consequence, on a selected contract, of a machine-verified theorem.

In this configuration, ArchSig's design principle of ``staying silent
about what cannot be spoken'' works as a safety device. Agents read
outputs literally and act on them literally, so a checker that warns
without grounds makes the generative loop diverge. Only a checker that
speaks exactly what can be spoken from the given contract can be
placed as a gate of a generative system. Economically as well, the
computational cost of deterministic verification is small compared to
probabilistic generation, and the faster generation becomes, the
higher the value of computation that rejects mistakes before
execution.

Human work does not disappear in this configuration; it moves. Which
vocabulary is admitted, which equations are required, what counts as
PASS --- the selection of the LawPolicy remains on the human side as
the decision of what deserves protection, and is indeed purified into
it. The diagnosis staircase of Section~\ref{sec:7} (nonzero residual
$\to$ BLOCKED $\to$ prior verification of the repair $\to$ PASS) is
the minimal working example of this configuration.

\subsection{The Rising Sea research program}\label{sec:96}

The research program as a whole is the following chain:
\begin{center}
\small
put architectural facts into words
$\to$ construct AAT geometry from Atoms and equation systems
$\to$ formalize the mathematical claims in Lean
$\to$ measure selected finite instances
$\to$ SFT connects the geometry to evolution dynamics
$\to$ compute trajectories, reachable futures, feedback, governance
\end{center}

SAGA is the first instance traversing the first three stages of this
chain --- mathematics, formalization, measurement --- in a single
theorem. The provenance discipline this paper has fixed (classifying
the kinds of claims and connecting each claim to primary evidence)
will be used in the same form on all subsequent peaks. When the water
level of the theory has risen high enough, users see only the surface.
An engineer, having learned neither categories nor sheaves nor
cohomology, says ``analyze this'', receives only the word
\emph{obstruction}, fixes it, and moves on. That quietness is the goal
of this program.
